% sec-01: outline item A.  Figures 1 (mermaid), 2 (d2), 3 (graphviz).
\section{The Novel-Performer Case}

\subsection{Three Clauses, One Applicant}

\emph{Science: A New Golden Age} contains three clauses under which a Phase 1
pancreatic cancer trial could plausibly be funded, and they are not
interchangeable \cite{kratsios2026sciencegoldenage}.

The first is the individual-scientist language of Chapter I, which asks the
roughly \$200 billion annual federal research and development portfolio to stop
deferring to ``the same incumbents that consume the funding''. The ten prior
applications were written under it \cite{tenapps2026}.

The second is the chapter section headed \textsc{novel performers}, which
describes ``a vast middle ground of mid-scale science'', problems ``requiring
tens of millions of dollars, coordinated teams of ten to a hundred people, and
timelines of half a decade''.

The third is the Small Business Innovation Research clause, and it is the only
one of the three that appears in more than one chapter. The Chapter I summary of
recommendations asks agencies to ``focus Small Business Innovation Research
(SBIR) and Small Business Technology Transfer (STTR) programs to build strategic
capabilities''. Chapter III states that ``SBIR and STTR programs can be deployed
strategically to advance new scientific and technological capabilities, coupling
federally-seeded companies with the scientific enterprise''. Chapter IV states
that ``Programs like SBIR open doors for technician-founded ventures, directing
resources toward ideas that do not always start with a dissertation''.

\begin{appfloat}[!tb]
\begin{appfig}[mermaid-type / flowchart]
% Single decision spine at y = 0, diamonds on a 3.15cm pitch so 4.9mm of clear
% canvas separates adjacent diamond tips.  Three drop-outs on a second band at
% y = -2.9.  No edge crosses another: the spine is horizontal and the three
% exits are near-vertical.
\node[mmgoal,text width=26mm] (src) at (0,0)
  {Science: A New Golden Age\\three candidate clauses};
\foreach \i/\x/\lab in {1/3.3/{a firm?}, 2/6.45/{ask 10 to 50M?},
  3/9.60/{F and A base?}, 4/12.75/{team of 10 to 100?}}{
  \node[mmdec,text width=14mm,aspect=2.2] (t\i) at (\x,0) {\lab};
  \node[font=\tiny\sffamily\bfseries,text=protoblue,anchor=south]
       at ([yshift=0.6mm]t\i.north) {T\i};}
\node[mmgray,text width=28mm] (f1) at (2.6,-2.9)
  {Clause 1 ends\\the applicant is now an entity, not a person};
\node[mmgray,text width=30mm] (f2) at (7.5,-2.9)
  {Clause 2 ends\\written for tens of millions; would also fail T4};
\node[mmgoal,text width=30mm] (pass) at (12.75,-2.9)
  {Clause 3 survives\\NIH SEED SBIR, 306K then 1.3M};
\draw[mmedgeb] (src) -- (t1);
\draw[mmedgeb] (t1) -- node[mmlabel] {yes} (t2);
\draw[mmedgeb] (t2) -- node[mmlabel] {no, 1.6M} (t3);
\draw[mmedgeb] (t3) -- node[mmlabel] {none} (t4);
\draw[mmedged] (t1.south) -- node[mmlabel,pos=0.55] {yes} (f1.north);
\draw[mmedged] (t2.south) -- node[mmlabel,pos=0.55] {no} (f2.north);
\draw[mmedgeb,line width=1pt] (t4.south) -- node[mmlabel,pos=0.55] {no, 2.6 FTE} (pass.north);
\node[font=\tiny\itshape,text=protogray,anchor=north west,text width=132mm]
  at (0,-4.35) {The two clauses that fail, fail on size and on legal form, not
  on merit or on subject matter. T3 is the only test this company passes by
  lacking something rather than by having it.};
\end{appfig}
\vspace{-0.65cm}
\figcaption{Three clauses of one report, four eligibility tests, and the only\\
clause a San Diego firm with no indirect-cost base survives. The\\
two that fail, fail on size, not on merit or on subject matter.}
\end{appfloat}

\begin{apptable}
\begin{tabularx}{\textwidth}{>{\raggedright\arraybackslash}p{0.7cm} >{\raggedright\arraybackslash}p{3.5cm} >{\raggedright\arraybackslash}X >{\raggedright\arraybackslash}p{3.0cm}}
\headrow \thead{Test} & \thead{Question} & \thead{Basis in the report} & \thead{ChemicalQDevice, August 2026}\\
\midrule
T1 & Is the applicant a firm rather than a person? & Chapter IV, SBIR opens doors for technician-founded ventures & Yes, a California corporation \\
T2 & Is the ask ten to fifty million? & Chapter II, novel performers is written for tens of millions & No, \$1,606,000 over 33 months \\
T3 & Is there an indirect-cost base to recover? & Chapter II, deference to the incumbents that consume the funding & No negotiated rate agreement exists \\
T4 & Is the team ten to a hundred people? & Chapter II, coordinated teams of ten to a hundred & No, 2.6 FTE at full Phase II staffing \\
\end{tabularx}
\end{apptable}
\tabcap{The four tests, the sentence in the report that each is drawn from, and
this company's answer. Clause 1 fails T1, clause 2 fails T2 and T4, and clause 3
is the only one that survives all four.}

\subsection{Why the Pioneer Language Does Not Fit an Entity}

What changed between the ten applications and this one is the applicant. A
person-based mechanism funds a named scientist and follows that person's
judgement for five to seven years. ChemicalQDevice is a California corporation
with one officer, and a corporation is not a subject that mechanism has a way to
fund. Nothing about the science changed; the legal form of the applicant did.

That is a real loss and it is worth naming here rather than in a footnote. The
individual-scientist framing is what made all ten prior applications distinctive,
and it is the framing the report itself rewards most explicitly. Abandoning it
is the first of four costs the conversion carries, and all four are set out in
\S9 without mitigation.

\subsection{Why the Novel-Performers Chapter Does Not Fit Either}

The novel-performers chapter is precise about scale, and the precision is what
excludes this company. It describes problems requiring tens of millions of
dollars, teams of ten to a hundred people, and timelines of half a decade, and
it names focused research organizations and the NSF TIP X-Labs as the vehicles
built for them. Set against that: \$1,606,000, 2.6 full-time equivalents, and 33
months.

The chapter describes a real gap in the federal portfolio. This company does not
fill it, and claiming otherwise would be exactly the kind of overreach the
report is written against.

\begin{appfloat}[!tb]
\begin{appfig}[d2-type / true grid]
% One true grid.  All five score columns take the identical 20mm width on a
% 2.1cm pitch, so the header row cannot drift out of register with the body.
% Row pitch 0.78cm against a 0.72cm cell height leaves a 0.6mm rule gap.
\foreach \c/\x/\lab in {0/1.75/{Activity}, 1/4.55/{University}, 2/6.65/{Corporate lab},
  3/8.75/{Federal lab}, 4/10.85/{New institutions}}{
  \ifnum\c=0
    \node[d2cellh,text width=33mm,minimum height=7.2mm] (h\c) at (\x,0) {\lab};
  \else
    \node[d2cellh,minimum width=20mm,text width=17mm,minimum height=7.2mm] (h\c) at (\x,0) {\lab};
  \fi}
\node[d2cellk,minimum width=20mm,text width=17mm,minimum height=7.2mm,
      font=\tiny\sffamily\bfseries] (h5) at (12.95,0) {SBIR firm 2.6 FTE};
\foreach \r/\y/\lab in {%
  1/-0.78/{Curiosity driven, investigator led},
  2/-1.56/{Larger scale, engineering intensive},
  3/-2.34/{Long horizon platform and tools},
  4/-3.12/{Public goods data and infrastructure},
  5/-3.90/{Mission driven, public good science},
  6/-4.68/{Proprietary product development},
  7/-5.46/{Workforce training and apprenticeship}}{
  \node[d2celll,text width=33mm,minimum height=7.2mm] at (1.75,\y) {\lab};}
\foreach \x/\y/\sty/\lab in {%
  4.55/-0.78/d2cell/{Well suited}, 6.65/-0.78/d2cellg/{Less suited},
  8.75/-0.78/d2cell/{Partly suited}, 10.85/-0.78/d2cell/{By design},
  4.55/-1.56/d2cellg/{Less suited}, 6.65/-1.56/d2cell/{Partly suited},
  8.75/-1.56/d2cell/{Partly suited}, 10.85/-1.56/d2cell/{By design},
  4.55/-2.34/d2cellg/{Less suited}, 6.65/-2.34/d2cell/{Partly suited},
  8.75/-2.34/d2cell/{Well suited}, 10.85/-2.34/d2cell/{By design},
  4.55/-3.12/d2cell/{Partly suited}, 6.65/-3.12/d2cellg/{Less suited},
  8.75/-3.12/d2cell/{Partly suited}, 10.85/-3.12/d2cell/{By design},
  4.55/-3.90/d2cell/{Partly suited}, 6.65/-3.90/d2cellg/{Less suited},
  8.75/-3.90/d2cell/{Partly suited}, 10.85/-3.90/d2cell/{By design},
  4.55/-4.68/d2cellg/{Less suited}, 6.65/-4.68/d2cell/{Well suited},
  8.75/-4.68/d2cellg/{Less suited}, 10.85/-4.68/d2cell/{By design},
  4.55/-5.46/d2cell/{Well suited}, 6.65/-5.46/d2cell/{Partly suited},
  8.75/-5.46/d2cellg/{Less suited}, 10.85/-5.46/d2cell/{By design}}{
  \node[\sty,minimum width=20mm,text width=17mm,minimum height=7.2mm] at (\x,\y) {\lab};}
\foreach \y/\sty/\lab in {%
  -0.78/d2soft/{Partly suited}, -1.56/d2gray2/{Less suited},
  -2.34/d2soft/{Partly suited}, -3.12/d2key/{Well suited},
  -3.90/d2key/{Well suited}, -4.68/d2soft/{Partly suited},
  -5.46/d2gray2/{Less suited}}{
  \node[\sty,minimum width=20mm,text width=17mm,minimum height=7.2mm] at (12.95,\y) {\lab};}
\draw[protoblue,line width=0.9pt] (11.90,0.42) -- (11.90,-5.85);
\legkey{0.0}{-6.45}{protowhite}{Well or partly suited, the report's own scoring}
\legkey{6.6}{-6.45}{pagrayl}{Less suited}
\legkey{9.6}{-6.45}{protoblue}{Well suited, this company}
\node[font=\tiny\itshape,text=protogray,anchor=north west,text width=132mm]
  at (0,-6.85) {The rule at the right marks where the report's own columns end
  and the column this paper adds begins. The four left-hand columns are the
  report's scoring restated in words, because its plus, tilde and cross glyphs
  are not legible at this measure.};
\end{appfig}
\vspace{-0.65cm}
\figcaption{The report's own seven activities, its four institutional forms,\\
and a fifth column it does not carry. An SBIR firm of 2.6 people\\
is well suited to two of the seven and unsuited to two others.}
\end{appfloat}

Two well suited, three partly suited, two less suited. The two well-suited rows
are public-goods data and infrastructure, and mission-driven public-good
science, and both are well suited for the same structural reason: every one of
the twelve milestone artifacts in \S5 is deposited publicly by plan, and the
programme has one mission with a stop condition and a closure budget.

A firm claiming to be well suited to all seven would be claiming to be a
university, a corporate laboratory, a federal laboratory and a focused research
organization at once. The report's own table exists to make that claim
impossible to state.

\subsection{The Clause That Fits, and What It Is Worth}

ChemicalQDevice has no negotiated indirect-cost rate agreement, no facilities to
recover, and no administrative layer to fund. That is a structural fact about a
2.6 person firm rather than a concession, and it has a price a program officer
with a fixed appropriation can read directly.

\begin{appfloat}[!tb]
\begin{appfig}[graphviz-type / record nodes]
% Four record nodes, each a stack of ruled cells sharing borders.  Label column
% 21mm, value column 12mm, so the four totals sit on one vertical line at
% x offset 33mm inside every record.  Record pitch 3.75cm against a 3.3cm
% record width leaves 4.5mm of clear canvas between records.
\foreach \rx/\hdr/\hsty in {0/{Direct work, identical}/gvcellh,
  3.75/{Route A, university}/gvcellh, 7.50/{Route B, full allowance}/gvcellh,
  11.25/{Route C, this plan}/gvcellh}{
  \node[\hsty,minimum width=34mm,text width=31mm,minimum height=6.4mm,anchor=north west]
       at (\rx,0.64) {\hdr};}
\foreach \i/\a/\b in {1/{Personnel, 2.6 FTE}/{\$842,000}, 2/{Site clinical conduct}/{\$288,000},
  3/{Compute and rig}/{\$96,000}, 4/{Monitoring and DSMB}/{\$86,000},
  5/{Consultants, travel}/{\$84,000}}{
  \node[gvcells,minimum width=21mm,text width=18mm,minimum height=5.0mm,anchor=north west,align=left]
       at (0,{-5.0*(\i-1)/10}) {\a};
  \node[gvcells,minimum width=13mm,text width=10.6mm,minimum height=5.0mm,anchor=north west]
       at (2.1,{-5.0*(\i-1)/10}) {\b};}
\node[gvcellh,minimum width=21mm,text width=18mm,minimum height=5.4mm,anchor=north west,align=left]
     at (0,-2.5) {Total direct};
\node[gvcellh,minimum width=13mm,text width=10.6mm,minimum height=5.4mm,anchor=north west]
     at (2.1,-2.5) {\$1,396,000};
\foreach \i/\a/\b in {1/{Direct}/{\$1,396,000}, 2/{Less equipment}/{-\$96,000},
  3/{MTDC base}/{\$1,300,000}, 4/{F and A, 57 pct}/{\$741,000}, 5/{Fee}/{none}}{
  \node[gvcellg,minimum width=21mm,text width=18mm,minimum height=5.0mm,anchor=north west,align=left]
       at (3.75,{-5.0*(\i-1)/10}) {\a};
  \node[gvcellg,minimum width=13mm,text width=10.6mm,minimum height=5.0mm,anchor=north west]
       at (5.85,{-5.0*(\i-1)/10}) {\b};}
\node[gvcell,fill=pagraym,minimum width=21mm,text width=18mm,minimum height=5.4mm,anchor=north west,align=left]
     at (3.75,-2.5) {Total to funder};
\node[gvcell,fill=pagraym,minimum width=13mm,text width=10.6mm,minimum height=5.4mm,anchor=north west]
     at (5.85,-2.5) {\$2,137,000};
\foreach \i/\a/\b in {1/{Direct}/{\$1,396,000}, 2/{Indirect, 40 pct}/{\$558,000},
  3/{Subtotal}/{\$1,954,000}, 4/{Fee, 7 pct}/{\$137,000}, 5/{}/{}}{
  \node[gvcellg,minimum width=21mm,text width=18mm,minimum height=5.0mm,anchor=north west,align=left]
       at (7.50,{-5.0*(\i-1)/10}) {\a};
  \node[gvcellg,minimum width=13mm,text width=10.6mm,minimum height=5.0mm,anchor=north west]
       at (9.60,{-5.0*(\i-1)/10}) {\b};}
\node[gvcell,fill=pagraym,minimum width=21mm,text width=18mm,minimum height=5.4mm,anchor=north west,align=left]
     at (7.50,-2.5) {Total to funder};
\node[gvcell,fill=pagraym,minimum width=13mm,text width=10.6mm,minimum height=5.4mm,anchor=north west]
     at (9.60,-2.5) {\$2,091,000};
\foreach \i/\a/\b in {1/{Direct}/{\$1,396,000}, 2/{Indirect, 7.5 pct}/{\$105,000},
  3/{Subtotal}/{\$1,501,000}, 4/{Fee, 7 pct}/{\$105,000}, 5/{}/{}}{
  \node[gvcells,minimum width=21mm,text width=18mm,minimum height=5.0mm,anchor=north west,align=left]
       at (11.25,{-5.0*(\i-1)/10}) {\a};
  \node[gvcells,minimum width=13mm,text width=10.6mm,minimum height=5.0mm,anchor=north west]
       at (13.35,{-5.0*(\i-1)/10}) {\b};}
\node[gvcellh,minimum width=21mm,text width=18mm,minimum height=5.4mm,anchor=north west,align=left]
     at (11.25,-2.5) {Total to funder};
\node[gvcellh,minimum width=13mm,text width=10.6mm,minimum height=5.4mm,anchor=north west]
     at (13.35,-2.5) {\$1,606,000};
\draw[gvedgeb] (1.65,-3.28) to[bend right=26]
  node[font=\tiny,fill=protowhite,inner sep=1.5pt,pos=0.5] {premium \$531,000, 33.1 percent} (5.55,-3.28);
\draw[gvedgeb] (9.30,-3.28) to[bend left=26]
  node[font=\tiny,fill=protowhite,inner sep=1.5pt,pos=0.5] {premium \$485,000, 30.2 percent} (13.05,-3.28);
\begin{scope}[shift={(0,-5.7)}]
  \axisx{0}{4.4}{0}
  \vbarcol{0.75}{2.14}{pagraym}{Route A}{2,137}
  \vbarcol{2.20}{2.09}{pagraym}{Route B}{2,091}
  \vbarcol{3.65}{1.61}{protoblue}{Route C}{1,606}
  \node[font=\tiny\sffamily,text=protogray,anchor=west] at (4.6,1.2)
    {Bar height is thousands of dollars to the funder for one identical
    \$1,396,000 of direct work.};
\end{scope}
\node[font=\tiny\itshape,text=protogray,anchor=north west,text width=132mm]
  at (0,-6.35) {Every column begins at the same \$1,396,000 and every total is
  the sum of its own column. The company may claim route B and does not; the
  difference between B and C is \$485,000 that is never requested.};
\end{appfig}
\vspace{-0.65cm}
\figcaption{The same 1,396,000 of direct work priced three ways. A university\\
F and A rate at 57 percent adds 741,000, the full SBIR allowance\\
adds 695,000, and the rate this plan claims adds 210,000 in all.}
\end{appfloat}

\begin{apptable}
\begin{tabularx}{\textwidth}{>{\raggedright\arraybackslash}X >{\raggedright\arraybackslash}p{2.6cm} >{\raggedright\arraybackslash}p{2.9cm} >{\raggedright\arraybackslash}p{2.6cm}}
\headrow \thead{Line} & \thead{Route A, university} & \thead{Route B, full allowance} & \thead{Route C, this plan}\\
\midrule
Direct costs & \$1,396,000 & \$1,396,000 & \$1,396,000 \\
Equipment removed from base & -\$96,000 & not applicable & not applicable \\
Overhead base & \$1,300,000 MTDC & \$1,396,000 direct & \$1,396,000 direct \\
Overhead rate & 57 percent & 40 percent & 7.5 percent \\
Overhead & \$741,000 & \$558,000 & \$105,000 \\
Fee at 7 percent & none & \$137,000 & \$105,000 \\
Total to the funder & \$2,137,000 & \$2,091,000 & \$1,606,000 \\
Premium over route C & \$531,000, 33.1 pct & \$485,000, 30.2 pct & zero \\
\end{tabularx}
\end{apptable}
\tabcap{The same direct work under three overhead regimes. The 57 percent rate
is a typical negotiated university rate; the 40 percent is the SBIR allowance
available without a rate agreement; the 7.5 percent is what this firm documents.}

The 7.5 percent rate covers insurance, audit and accounting, and nothing else,
because there is nothing else to cover \cite{nihsbirseed}. The company could
claim the 40 percent allowance without documentation and does not
\cite{ecfr2026uniformguidance}. For a program officer holding a fixed
appropriation the consequence is arithmetic rather than rhetorical: the same
trial, run by the same person, costs \$531,000 less than it would through an
institution and \$485,000 less than this mechanism permits.
